% =============================================================================
% python-KernSmooth Technical Report
% Main source file: docs/report/main.tex
%
% Compilation (from the project root):
%   pdflatex docs/report/main.tex
%   pdflatex docs/report/main.tex   (second pass for cross-references)
%
% All \lstinputlisting paths are relative to this file's directory (docs/report/).
% Relative paths of the form ../../.claude/ therefore resolve to .claude/ under
% the project root, which is the correct location of the agent/command specs.
% =============================================================================

\documentclass[12pt,a4paper]{article}

% --- Encoding and fonts ---
\usepackage[T1]{fontenc}
\usepackage[utf8]{inputenc}
\usepackage{lmodern}
\usepackage{microtype}

% --- Page layout ---
\usepackage[top=1.2in,bottom=1.2in,left=1.1in,right=1.1in]{geometry}

% --- Math ---
\usepackage{amsmath}
\usepackage{amssymb}

% --- Tables ---
\usepackage{booktabs}
\usepackage{array}
\usepackage{longtable}

% --- Color ---
\usepackage{xcolor}
\definecolor{bgcode}{RGB}{248,248,248}
\definecolor{fgborder}{RGB}{200,200,200}
\definecolor{appbg}{RGB}{241,246,255}
\definecolor{appborder}{RGB}{110,145,210}
\definecolor{kwblue}{RGB}{0,40,170}
\definecolor{strred}{RGB}{160,0,0}
\definecolor{comgray}{RGB}{110,110,110}
\definecolor{linkblue}{RGB}{10,80,160}

% --- Code listings ---
\usepackage{listings}

%% Generic monospace style (shell commands, build scripts, etc.)
\lstdefinestyle{shell}{
  basicstyle=\ttfamily\small,
  backgroundcolor=\color{bgcode},
  frame=single,
  rulecolor=\color{fgborder},
  breaklines=true,
  breakatwhitespace=true,
  columns=flexible,
  keepspaces=true,
  showstringspaces=false,
  xleftmargin=6pt,
  xrightmargin=6pt,
  framexleftmargin=4pt,
}

%% Python syntax-highlighted style
\lstdefinestyle{python}{
  language=Python,
  basicstyle=\ttfamily\small,
  backgroundcolor=\color{bgcode},
  frame=single,
  rulecolor=\color{fgborder},
  breaklines=true,
  columns=flexible,
  keepspaces=true,
  showstringspaces=false,
  xleftmargin=6pt,
  xrightmargin=6pt,
  framexleftmargin=4pt,
  keywordstyle=\color{kwblue}\bfseries,
  stringstyle=\color{strred},
  commentstyle=\color{comgray}\itshape,
}

%% Language definition for Markdown files (used in the appendix)
%  No keyword highlighting - we want verbatim reproduction with a framed box.
\lstdefinelanguage{Markdown}{
  basicstyle=\ttfamily\scriptsize,
  breaklines=true,
  breakatwhitespace=false,
  columns=fullflexible,
  keepspaces=true,
  tabsize=2,
  showstringspaces=false,
}

%% Appendix-file listing style
\lstdefinestyle{mdfile}{
  language=Markdown,
  backgroundcolor=\color{appbg},
  frame=single,
  rulecolor=\color{appborder},
  framesep=5pt,
  xleftmargin=4pt,
  xrightmargin=4pt,
  aboveskip=6pt,
  belowskip=6pt,
}

% --- Headers and footers ---
\usepackage{fancyhdr}
\setlength{\headheight}{13.59999pt}
\addtolength{\topmargin}{-1.59999pt}
\pagestyle{fancy}
\fancyhf{}
\renewcommand{\headrulewidth}{0.4pt}
\fancyhead[L]{\small\itshape python-KernSmooth}
\fancyhead[R]{\small\itshape \nouppercase{\leftmark}}
\fancyfoot[C]{\small\thepage}
\renewcommand{\sectionmark}[1]{\markboth{#1}{}}

% --- Appendix package ---
\usepackage[toc,page]{appendix}

% --- Paragraph spacing ---
\usepackage{parskip}
\setlength{\parindent}{0pt}

% --- List formatting ---
\usepackage{enumitem}
\setlist[itemize]{noitemsep,topsep=4pt}
\setlist[enumerate]{noitemsep,topsep=4pt}
\setlist[description]{noitemsep,topsep=4pt}

% --- Hyperref (load last) ---
\usepackage{hyperref}
\hypersetup{
  colorlinks  = true,
  linkcolor   = linkblue,
  urlcolor    = linkblue,
  citecolor   = {green!45!black},
  pdftitle    = {python-KernSmooth: A Technical Report},
  pdfauthor   = {Yufei Cai},
  pdfsubject  = {Porting KernSmooth from R to Python},
  pdfkeywords = {KernSmooth, R, Python, kernel smoothing, f2py, Fortran},
}

% =============================================================================
% Custom commands
% =============================================================================

%% \appendixentry{Display Title}{relative/path/to/file.md}
%%
%% Inserts a labelled subsubsection in the appendix, adds it to the TOC,
%% and typesets the file contents verbatim using the mdfile listing style.
%% The label is app:<Display Title> (with spaces and special chars as given;
%% callers should pass simple hyphenated names for safety).
\newcommand{\appendixentry}[2]{%
  \phantomsection%
  \subsubsection*{\texttt{#1}}%
  \addcontentsline{toc}{subsubsection}{\texttt{#1}}%
  \label{appentry:#1}%
  \lstinputlisting[style=mdfile]{#2}%
  \vspace{0.5em}%
}

% =============================================================================
\begin{document}
% =============================================================================

% --- Title page ---
\begin{titlepage}
  \centering
  \vspace*{3.5cm}
  {\Huge\bfseries python-KernSmooth\par}
  \vspace{0.8cm}
  {\LARGE A Technical Report on Porting the\\[0.4em]
   \texttt{KernSmooth} R Package to Python\par}
  \vspace{1.4cm}
  {\Large Yufei Cai\par}
  \vspace{0.5cm}
  {\normalsize\ttfamily /groups/jli9/Yufei/python-KernSmooth\par}
  \vfill
  {\normalsize\textit{Status:~Work in Progress}\par}
  \vspace{0.3cm}
  {\normalsize Last updated:~2026-05-20\par}
\end{titlepage}

\tableofcontents
\clearpage

% =============================================================================
\section*{Abstract}
\addcontentsline{toc}{section}{Abstract}
% =============================================================================

This report documents the ongoing work to port the R package \texttt{KernSmooth} (version~2.23-26, released 2024-12-10) into an installable Python package named \texttt{r2py\_kernsmooth}. \texttt{KernSmooth} is a \emph{recommended} package shipped with every standard R installation; it implements kernel smoothing methods from Wand \& Jones~\cite{wandjones1995} and has no runtime dependencies on any external CRAN package.

The conversion strategy retains the original Fortran~77 computational core verbatim, compiles it into a Python C-extension via \texttt{numpy.f2py}, and replaces the R wrapper layer with semantically equivalent Python functions. The project proceeds in five phases: (1)~static architectural analysis of the R package, (2)~Python build-system establishment, (3)~language-dependency cataloguing and per-dependency conversion guide generation, (4)~automated R-to-Python function conversion and package assembly, and (5)~code-quality audit and namespace refinement. All five phases have been executed. The resulting package installs under Python~3.14 (CPython) in a \texttt{conda} environment on a Red Hat Enterprise Linux~9 HPC node and produces numerically correct output.

% =============================================================================
\section{Introduction}
% =============================================================================

\subsection{Background and Motivation}

The R package \texttt{KernSmooth} implements the statistical methods described in
\begin{quote}
  M.~P.~Wand and M.~C.~Jones (1995).\ \textit{Kernel Smoothing}.
  Chapman and Hall, London.
\end{quote}
It is a \emph{recommended} package --- one of the roughly fifteen packages that ship as part of the base R distribution and are maintained by the R Core Team. The package was authored by M.~P.~Wand, contributed LINPACK Fortran routines are attributed to Cleve Moler, and the current R-language port and maintenance is by Brian Ripley. It is released under an ``Unlimited'' license.

The package provides seven public functions covering kernel density estimation, plug-in bandwidth and bin-width selection, and local polynomial regression. Because it delegates all numerically intensive loops to compiled Fortran~77 code, it is both correct and performant.

The motivation for a Python port is to make these well-tested implementations available to Python-based scientific workflows without requiring an R runtime. A port that re-uses the \emph{original Fortran routines} offers a stronger correctness guarantee than a ground-up Python reimplementation: numerical differences between the Python and R packages can then originate only in the R-to-Python wrapper translation, which is far easier to audit systematically.

The broader scientific motivation, experimental design, and validation plan for this project are described in two planning documents that are not reproduced in this report:
\begin{itemize}
  \item \texttt{docs/methodology\_and\_plan\_Jun.pdf} --- the primary methodology and project plan.
  \item \texttt{docs/methodology\_addendum\_20260511.pdf} --- an addendum dated 2026-05-11.
\end{itemize}
The original R package reference manual is archived at \texttt{docs/KernSmooth.pdf}. The contents of these PDF files were not available during the writing of this report and are therefore not summarised here.

\subsection{Repository Structure}

The working directory \texttt{python-KernSmooth/} is the working tree of the \texttt{python-KernSmooth} repository hosted at \texttt{github.com/caiyufei8/python-KernSmooth}. The \texttt{r2py\_kernsmooth/} subdirectory is additionally managed as a git subtree linked to a separate repository at \texttt{github.com/caiyufei8/r2py\_kernsmooth}, allowing the Python package to be versioned and published to PyPI independently of the broader project. The cluster scripts \texttt{git\_pull.sh} and \texttt{git\_push.sh} at the project root encapsulate the subtree synchronisation commands required to keep both remotes up to date from the HPC node.

\begin{center}
\begin{tabular}{p{0.42\textwidth}p{0.50\textwidth}}
\toprule
Path & Contents \\
\midrule
\texttt{KernSmooth/} &
  Unmodified R package source (v2.23-26). \\
\texttt{r2py\_kernsmooth/} &
  Git subtree linked to \texttt{github.com/caiyufei8/r2py\_kernsmooth}; the installable Python package. Internal structure detailed in the table below. \\
\texttt{structural\_analysis/} &
  Dependency-graph visualisation artefacts (\texttt{dependency\_graph.\{py,html,pdf,png\}}), the function-level CSV (\texttt{dependency\_levels.csv}), and the JSON analysis outputs per R file. \\
\texttt{language\_dependency\_analysis/} &
  Per-file CSV dependency tables, the CSV merge script, the combined table, and the 46 per-dependency conversion guides in \texttt{conversion\_guides/}. \\
\texttt{conversion\_results/} &
  Per-function JSON translation artefacts produced by the conversion pipeline. \\
\texttt{docs/} &
  Architecture document, planning PDFs, per-phase summaries in \texttt{daily\_summaries/}, and this report in \texttt{report/}. \\
\texttt{.claude/agents/} &
  Sub-agent specification Markdown files. \\
\texttt{.claude/commands/} &
  Skill (slash command) specification Markdown files. \\
\texttt{git\_pull.sh} &
  Cluster batch script; pulls from the \texttt{python-KernSmooth} main remote and merges changes from the \texttt{r2py\_kernsmooth} subtree remote into \texttt{r2py\_kernsmooth/}. \\
\texttt{git\_push.sh} &
  Cluster batch script; pushes to the \texttt{python-KernSmooth} main remote and propagates \texttt{r2py\_kernsmooth/} changes to the \texttt{r2py\_kernsmooth} subtree remote via \texttt{git subtree push}. \\
\bottomrule
\end{tabular}
\end{center}

The internal structure of \texttt{r2py\_kernsmooth/} is as follows.

\begin{center}
\begin{tabular}{p{0.46\textwidth}p{0.46\textwidth}}
\toprule
Path (relative to \texttt{r2py\_kernsmooth/}) & Contents \\
\midrule
\texttt{r2py\_kernsmooth/\_\_init\_\_.py} &
  The main Python module; defines the public API (\texttt{\_\_all\_\_}) and all wrapper functions. \\
\texttt{meson.build} &
  Meson build definition: f2py wrapper generation, NumPy include discovery, BLAS detection, and extension module compilation. \\
\texttt{pyproject.toml} &
  Package metadata, build-system declaration (\texttt{meson-python}), runtime dependencies, and cibuildwheel platform configuration. \\
\texttt{src/} &
  Thirteen Fortran~77 source files: eleven from the original \texttt{KernSmooth} package and two LINPACK routines sourced from Netlib. \\
\texttt{tests/} &
  Test scripts for validating the package against the R \texttt{KernSmooth} package. \\
\texttt{pip\_install.sh} &
  Cluster batch script for local development installation via \texttt{pip install --no-build-isolation}. \\
\texttt{.github/workflows/python-publish.yml} &
  GitHub Actions workflow; builds binary wheels for five platforms and publishes them to PyPI on every release via OIDC trusted publishing. \\
\bottomrule
\end{tabular}
\end{center}

\subsection{Automated Workflow Infrastructure}

This project employs \textbf{Claude Code} (Anthropic) as an AI-assisted development environment. Recurring multi-step tasks are encoded as \emph{skills} --- top-level slash commands whose specifications are stored as Markdown files under \texttt{.claude/commands/}. Skills orchestrate one or more \emph{sub-agents} whose specifications are stored under \texttt{.claude/agents/}. Both types of specification files use YAML front matter to declare a name and description, followed by a structured natural-language body describing execution steps and output schemas. Sub-agents are dispatched via the \texttt{@agent-name} notation from within a skill.

The full specifications for all agents and commands referenced in this report appear in Appendices~\ref{app:agents} and~\ref{app:commands}.

At the end of each phase, the \texttt{/summarize-research-report} skill was invoked:
\begin{lstlisting}[style=shell]
/summarize-research-report
\end{lstlisting}
This produced a structured Markdown summary of the phase's activities, which was saved to \texttt{docs/daily\_summaries/}. Those summaries are the primary source material for this report.

% =============================================================================
\section{Phase 1: Static Analysis of the KernSmooth R Package}
% =============================================================================

\subsection{Motivation}

Before any translation work could begin, a precise, machine-readable map of the R package's internal structure was required. The goals of this phase were:

\begin{enumerate}
  \item Classify every R function by dependency type: calls to R standard-library functions, calls to other functions defined within the same (or other self-written related) R source file, and calls through the Fortran foreign-function interface.
  \item Establish the public API surface and the complete internal call graph.
  \item Understand the Fortran integration layer and the symbol-registration convention it uses.
  \item Produce a comprehensive architecture document to serve as a reference throughout all subsequent phases.
\end{enumerate}

Without this foundation, later phases would need to infer architectural facts on a per-function basis, risking inconsistency and missed dependencies. The structural analysis JSON produced here also serves as a direct input to the language-dependency extraction skill in Phase~3 and to the conversion ordering logic in Phase~4.

\subsection{R Dependency Analysis}

The \texttt{/analyze-r-folder-dependencies} skill (Appendix~\ref{app:commands}, \S\ref{appentry:analyze-r-folder-dependencies}) was invoked against the R source directory:

\begin{lstlisting}[style=shell]
/analyze-r-folder-dependencies KernSmooth/R/
\end{lstlisting}

The skill performs a recursive scan of the target folder for \texttt{.R} and \texttt{.r} files, then dispatches the \texttt{@analyze-r-file-dependencies} sub-agent (Appendix~\ref{app:agents}, \S\ref{appentry:analyze-r-file-dependencies}) on each file. Because \texttt{KernSmooth} packages all R code into a single source file (\texttt{KernSmooth/R/all.R}), exactly one agent invocation was required.

The sub-agent parses the R file, identifies every user-defined function, and classifies each function call in its body into one of three mutually exclusive categories:

\begin{description}
  \item[\texttt{language\_dependencies}] Calls to R built-in functions or functions from installed packages, identified by namespace prefix or by not being locally defined (e.g.\ \texttt{fft}, \texttt{var}, \texttt{dnorm}, \texttt{stats::quantile}).
  \item[\texttt{internal\_dependencies}] Calls to other functions defined within the same R file.
  \item[\texttt{external\_dependencies}] Calls via \texttt{.Fortran()}, \texttt{.C()}, \texttt{.Call()}, or \texttt{.External()} to compiled native-code entry points.
\end{description}

The result is a JSON object keyed by function name, with each value containing the three dependency arrays, each listing unique function names for the dependencies. The output was written to \texttt{structural\_analysis/R/all.json}.

\subsection{Architecture Synthesis}

To produce a comprehensive architecture document, the following files were read in parallel and synthesised into \texttt{docs/architecture-analysis.md}:

\begin{center}
\begin{tabular}{p{0.46\textwidth}p{0.44\textwidth}}
\toprule
File & Purpose \\
\midrule
\texttt{KernSmooth/DESCRIPTION} &
  Package metadata, authorship, license, declared dependencies. \\
\texttt{KernSmooth/NAMESPACE} &
  Exported symbols; \texttt{useDynLib} registration directive. \\
\texttt{structural\_analysis/R/all.json} &
  Per-function dependency classification (produced above). \\
\texttt{structural\_analysis/dependency\_levels.csv} &
  Function call-depth hierarchy (pre-existing artefact). \\
\texttt{KernSmooth/man/bkde.Rd} &
  Density estimation API and algorithm description. \\
\texttt{KernSmooth/man/bkde2D.Rd} &
  Two-dimensional density estimation API. \\
\texttt{KernSmooth/man/bkfe.Rd} &
  Kernel functional estimator API. \\
\texttt{KernSmooth/man/dpih.Rd} &
  Histogram bin-width selector API. \\
\texttt{KernSmooth/man/dpik.Rd} &
  Kernel density bandwidth selector API. \\
\texttt{KernSmooth/man/dpill.Rd} &
  Local linear regression bandwidth selector API. \\
\texttt{KernSmooth/man/locpoly.Rd} &
  Local polynomial estimator API. \\
\bottomrule
\end{tabular}
\end{center}

A check for \texttt{KernSmooth/vignettes/} confirmed that the directory does not exist; the package has no vignette sources.

\subsection{Key Findings}

\subsubsection{Package Scope and External Dependencies}

\texttt{KernSmooth} has \emph{zero} runtime dependencies on contributed CRAN packages. It imports exactly five functions from the base \texttt{stats} package: \texttt{dbeta}, \texttt{dnorm}, \texttt{fft}, \texttt{quantile}, and \texttt{var}. \texttt{MASS} and \texttt{carData} appear in the \texttt{Suggests} field of \texttt{DESCRIPTION} but are referenced only in \texttt{man/} examples and play no role in the package's computational logic. This minimal dependency footprint simplifies the Python port considerably.

\subsubsection{Function Inventory}

\texttt{R/all.R} defines exactly 16 functions in a single file. Seven are exported via \texttt{NAMESPACE}; nine are internal (unexported).

\begin{center}
\begin{tabular}{llll}
\toprule
Function & Exported & Level & Fortran entry point \\
\midrule
\texttt{bkde}      & Yes & 0 & \texttt{F\_linbin} (via \texttt{linbin}) \\
\texttt{bkde2D}    & Yes & 0 & \texttt{F\_lbtwod} (via \texttt{linbin2D}) \\
\texttt{dpih}      & Yes & 0 & \texttt{F\_linbin} (via \texttt{linbin}) \\
\texttt{dpik}      & Yes & 0 & \texttt{F\_linbin} (via \texttt{linbin}) \\
\texttt{dpill}     & Yes & 0 & multiple (via six helpers) \\
\texttt{locpoly}   & Yes & 0 & \texttt{F\_locpol} and \texttt{F\_rlbin} \\
\texttt{bkfe}      & Yes & 1 & \texttt{F\_linbin} (via \texttt{linbin}) \\
\texttt{linbin}    & No  & 2 & \texttt{F\_linbin} \\
\texttt{rlbin}     & No  & 2 & \texttt{F\_rlbin} \\
\texttt{blkest}    & No  & 1 & \texttt{F\_blkest} \\
\texttt{cpblock}   & No  & 1 & \texttt{F\_cp} \\
\texttt{linbin2D}  & No  & 1 & \texttt{F\_lbtwod} \\
\texttt{sdiag}     & No  & 1 & \texttt{F\_sdiag} \\
\texttt{sstdiag}   & No  & 1 & \texttt{F\_sstdg} \\
\texttt{.onAttach} & No  & 0 & --- \\
\texttt{.onUnload} & No  & 0 & --- \\
\bottomrule
\end{tabular}
\end{center}

\texttt{bkfe} is notable for being both an \emph{exported} function and an \emph{internal dependency}: \texttt{dpih} and \texttt{dpik} both call it to evaluate the density functionals $\int f(x)\,f^{(r)}(x)\,\mathrm{d}x$ required for plug-in bandwidth selection.

\subsubsection{Three-Level Dependency Hierarchy}

The internal call graph has three levels, stored in \texttt{structural\_analysis/dependency\_levels.csv}:

\begin{description}
  \item[Level~0 (entry points):] Functions callable directly by the user: \texttt{bkde}, \texttt{bkde2D}, \texttt{dpih}, \texttt{dpik}, \texttt{dpill}, \texttt{.onAttach}, \texttt{.onUnload}.
  \item[Level~1 (mid-tier helpers):] Functions called by one or more level-0 functions: \texttt{bkfe}, \texttt{blkest}, \texttt{cpblock}, \texttt{linbin2D}, \texttt{locpoly}, \texttt{sdiag}, \texttt{sstdiag}.
  \item[Level~2 (binning primitives):] \texttt{linbin} and \texttt{rlbin}.
\end{description}

\texttt{linbin} (one-dimensional linear binning, backed by Fortran \texttt{F\_linbin}) is the single most central function: it is called by seven distinct callers (\texttt{bkde}, \texttt{bkfe}, \texttt{dpih}, \texttt{dpik}, \texttt{locpoly}, \texttt{sdiag}, \texttt{sstdiag}), meaning that every density estimation and functional estimation path passes through it.

\texttt{dpill} is the most internally complex exported function, orchestrating six internal helpers (\texttt{rlbin}, \texttt{cpblock}, \texttt{blkest}, \texttt{locpoly}, \texttt{sdiag}, \texttt{sstdiag}) to produce a single bandwidth estimate.

\subsubsection{Fortran Integration Convention}

The \texttt{NAMESPACE} file contains:

\begin{lstlisting}[style=shell]
useDynLib(KernSmooth, .registration = TRUE, .fixes = "F_")
\end{lstlisting}

This registers all Fortran entry points at package load time and applies the \texttt{F\_} prefix to their R-level names, preventing symbol conflicts. All eight Fortran-delegating R wrappers use this naming scheme. LINPACK-derived linear-algebra routines (attributed to Cleve Moler in \texttt{DESCRIPTION}) reside in \texttt{KernSmooth/src/d*} and underpin the local polynomial computations in \texttt{locpoly}, \texttt{sdiag}, and \texttt{sstdiag}.

\subsubsection{Architectural Design: Binned Approximation}

Every computationally non-trivial function in \texttt{KernSmooth} adopts the same two-phase strategy: (1)~\textbf{bin} the raw data onto an equally-spaced grid using \emph{linear binning}, which preserves the first moment of each data point within its bin interval (unlike simple histogram binning), and (2)~\textbf{convolve} the bin counts with a kernel weight vector, either via FFT (\texttt{bkde}, \texttt{bkde2D}, \texttt{bkfe}) or via the Fortran local-polynomial routines. This trades a small, controllable bias for a reduction in computational complexity from $O(n \cdot g)$ to $O(n + g \log g)$, where $n$ is the sample size and $g$ is the grid size.

\subsubsection{Note on Fortran Source Analysis}

A complete cross-language dependency map would also require analysing the compiled Fortran sources in \texttt{KernSmooth/src/}. The \texttt{/analyze-c-folder-dependencies} skill (Appendix~\ref{app:commands}) was identified as a candidate tool for this task (it handles both C and Fortran), but this analysis has not yet been executed. Due to the nature of \texttt{KernSmooth} package, the Fortran sources are relatively self-contained and do not call back into R, so the lack of a Fortran-level dependency map has not posed a significant obstacle to subsequent phases. The Fortran code is treated as a black box that is called by the R wrappers, so the focus has been on understanding the R-level structure and dependencies.

% =============================================================================
% =============================================================================
\section{Phase 2: Establishing the Python Build Infrastructure}
% =============================================================================

\subsection{Motivation}

The conversion strategy retains the original Fortran~77 sources verbatim and relies on \texttt{f2py} (part of NumPy) to parse them and generate a Python-callable C-extension. This approach provides a strong correctness guarantee: numerical differences between the Python and R packages can originate only in the R-to-Python wrapper translation, not in the underlying computational routines. Phase~2 addresses all build-system and packaging concerns before any wrapper code is written.

Two related objectives determine the scope of this phase. The first is local: the extension must compile and install correctly in the HPC development environment so that Python wrapper functions can be developed and tested. The second is distribution: once complete, the package must be available as pre-compiled binary wheels on PyPI for all major platforms, so that end users can install it with a plain \texttt{pip install r2py\_kernsmooth} without requiring a local Fortran compiler, BLAS library, or build toolchain. Because a Fortran extension produces a platform-specific binary, a separate wheel must be built for each combination of operating system, processor architecture, and CPython version.

\subsection{Build System Design}

The build backend is \texttt{meson-python}, declared in \texttt{r2py\_kernsmooth/pyproject.toml}:

\begin{lstlisting}[style=shell]
[build-system]
build-backend = "mesonpy"
requires = ["meson-python", "meson", "ninja", "numpy"]

[project]
name = "r2py_kernsmooth"
version = "0.1.0"
dependencies = ["numpy", "scipy"]
\end{lstlisting}

The \texttt{meson} and \texttt{ninja} entries in \texttt{requires} ensure that PEP~517 isolated builds can locate the Meson build system and the Ninja build tool without requiring them to be pre-installed in the user's environment. The complete build definition resides in \texttt{r2py\_kernsmooth/meson.build}.

\subsubsection{Fortran Sources}

Thirteen Fortran~77 fixed-form source files in \texttt{r2py\_kernsmooth/src/} are compiled into the extension:

\begin{center}
\begin{tabular}{p{0.36\textwidth}p{0.55\textwidth}}
\toprule
File & Origin \\
\midrule
\texttt{src/blkest.f}   & Original \texttt{KernSmooth} source \\
\texttt{src/cp.f}       & Original \texttt{KernSmooth} source \\
\texttt{src/dgedi.f}    & Original \texttt{KernSmooth} source (LINPACK LU) \\
\texttt{src/dgefa.f}    & Original \texttt{KernSmooth} source (LINPACK LU) \\
\texttt{src/dgesl.f}    & Original \texttt{KernSmooth} source (LINPACK LU) \\
\texttt{src/linbin.f}   & Original \texttt{KernSmooth} source \\
\texttt{src/linbin2D.f} & Original \texttt{KernSmooth} source \\
\texttt{src/locpoly.f}  & Original \texttt{KernSmooth} source \\
\texttt{src/rlbin.f}    & Original \texttt{KernSmooth} source \\
\texttt{src/sdiag.f}    & Original \texttt{KernSmooth} source \\
\texttt{src/sstdiag.f}  & Original \texttt{KernSmooth} source \\
\texttt{src/dqrdc.f}    & Added from Netlib (LINPACK QR decomposition, G.~W.~Stewart 1978) \\
\texttt{src/dqrsl.f}    & Added from Netlib (LINPACK QR solve, G.~W.~Stewart 1978) \\
\bottomrule
\end{tabular}
\end{center}

The two Netlib additions are necessary because \texttt{blkest.f} and \texttt{cp.f} call \texttt{dqrdc\_} and \texttt{dqrsl\_}, respectively. The original R package satisfies these symbols from R's own internal LINPACK copy (built into the R binary). OpenBLAS provides only BLAS levels~1--3 and does not include LINPACK QR routines. Sourcing them from Netlib is the standard remedy for this gap when porting R packages that depend on LINPACK.

The presence of LINPACK routines (both the original \texttt{dgedi/dgefa/dgesl} trio and the newly added \texttt{dqrdc/dqrsl}) was discovered iteratively during Phase~2: each link failure at \texttt{pip install} time revealed one missing symbol, which was then resolved by adding the corresponding Netlib file.

\subsubsection{f2py Wrapper Generation}

\texttt{f2py} (NumPy~2.4.3) is invoked as a Meson \texttt{custom\_target}:

\begin{lstlisting}[style=shell]
python -m numpy.f2py @INPUT@ -m _KernSmooth --lower
\end{lstlisting}

Because all sources are Fortran~77 fixed-form (\texttt{.f}), \texttt{f2py} generates \texttt{\_KernSmoothmodule.c} and \texttt{\_KernSmooth-f2pywrappers.f} (the fixed-form wrapper variant). An initial build template incorrectly named the wrapper file \texttt{\_KernSmooth-f2pywrappers2.f90} (the free-form Fortran~90 variant), which caused a build failure; the distinction was identified by inspecting the \texttt{f2py} output and corrected in the \texttt{meson.build} \texttt{outputs} list.

The \texttt{--lower} flag lowercases all Fortran symbol names in the generated C wrapper, which is required for standard GFortran symbol conventions.

\subsubsection{NumPy Include Directories}

The Meson configure step must locate NumPy's C header files at configure time (not at install time), since they are required to compile \texttt{fortranobject.c}, the f2py runtime support file. Two \texttt{run\_command} calls retrieve these directories:

\begin{lstlisting}[style=shell]
incdir_numpy = run_command(py,
  ['-c', 'import numpy; print(numpy.get_include())'],
  check: true).stdout().strip()
incdir_f2py = run_command(py,
  ['-c', 'import numpy.f2py; print(numpy.f2py.get_include())'],
  check: true).stdout().strip()
\end{lstlisting}

\texttt{fortranobject.c} (located in \texttt{incdir\_f2py}) is included directly in the extension module's source list.

\subsubsection{BLAS Discovery}
\label{sec:phase2-blas}

The HPC node (Red~Hat Enterprise Linux~9, GFortran~15.2.0, \texttt{ld.bfd~2.35.2-67}) provides OpenBLAS at \texttt{/usr/lib64/libopenblaso.so.0}, but the unversioned symlink (\texttt{libopenblaso.so}) normally installed by the \texttt{openblas-openmp-devel} package is absent. Meson's standard \texttt{find\_library} and \texttt{dependency('blas')} mechanisms cannot resolve a bare \texttt{.so.0} file. The following discovery chain in \texttt{meson.build} addresses this while remaining portable across all build platforms:

\begin{lstlisting}[style=shell]
blas_dep = dependency('blas', required: false)
if not blas_dep.found()
  blas_dep = fc.find_library('openblaso', required: false)
endif
if not blas_dep.found()
  blas_dep = fc.find_library('openblas', required: false)
endif
if not blas_dep.found()
  blas_dep = fc.find_library('blas', required: false)
endif
if not blas_dep.found()
  fs = import('fs')
  if fs.is_file('/usr/lib64/libopenblaso.so.0')
    blas_dep = declare_dependency(link_args:
      ['/usr/lib64/libopenblaso.so.0'])
  else
    error('BLAS not found. Install openblas-devel (Linux) '
        + 'or run: brew install openblas (macOS).')
  endif
endif
\end{lstlisting}

The first four stages cover systems where BLAS is installed conventionally: \texttt{dependency('blas')} queries \texttt{pkg-config}; the three \texttt{find\_library} calls search standard library directories for \texttt{openblaso}, \texttt{openblas}, and \texttt{blas} by name. The fifth and final stage is specific to the HPC cluster, where the \texttt{openblas-openmp-devel} package that would provide the unversioned symlink is absent but \texttt{/usr/lib64/libopenblaso.so.0} is present. Meson's \texttt{fs.is\_file()} guard verifies the file's existence before constructing the link argument, ensuring that on any platform where the hardcoded path is absent---such as macOS or a standard Ubuntu installation---the build fails at configure time with a clear, actionable error rather than proceeding silently to a linker failure.

\subsubsection{Package Structure}

The Python package exposes two components:

\begin{description}
  \item[\texttt{r2py\_kernsmooth/\_KernSmooth}] The compiled Fortran C-extension (private backend). All 11 original subroutines are accessible via \texttt{r2py\_kernsmooth.\_KernSmooth.\textit{name}}.
  \item[\texttt{r2py\_kernsmooth/\_\_init\_\_.py}] The Python wrapper layer that defines the public API.
\end{description}

An early design exposed the raw Fortran subroutines directly in the \texttt{r2py\_kernsmooth} namespace. This was identified as architecturally incorrect: the Fortran subroutines have a low-level calling convention (all arguments passed by reference, no input validation, no Python-idiomatic defaults) and should not be part of the public API. \texttt{\_\_init\_\_.py} was revised to contain only \texttt{from . import \_KernSmooth} and \texttt{\_\_all\_\_ = []}, deferring all public-API decisions to the wrapper functions to be written in Phase~4.

\subsection{Installation}

The package is installed into the \texttt{r-to-python} conda environment via:

\begin{lstlisting}[style=shell]
pip install --no-build-isolation .
\end{lstlisting}

The \texttt{--no-build-isolation} flag is required because the conda environment provides \texttt{numpy} and \texttt{meson-python} directly, and the Meson configure step calls \texttt{numpy.get\_include()} and \texttt{numpy.f2py.get\_include()} at configure time. With build isolation, pip would install \texttt{numpy} into an isolated build environment, producing a different path that may conflict with the runtime environment.

The full invocation, including the cluster's \texttt{module load gcc} and \texttt{conda activate r-to-python} steps, is captured in \texttt{r2py\_kernsmooth/pip\_install.sh}. This script is also suitable for submission as a cluster batch job (\texttt{\#\$ -pe smp 1}, \texttt{\#\$ -q long} directives are included).

After successful installation, \texttt{import r2py\_kernsmooth} and \texttt{import r2py\_kernsmooth.\_KernSmooth} both succeed without error.

\subsection{PyPI Distribution Infrastructure}

A package containing compiled Fortran code produces a platform-specific binary wheel that cannot be installed on a different operating system or processor architecture. Distributing pre-compiled wheels for all major platforms requires building them in CI on a representative of each target environment; the tool used for this purpose is \texttt{cibuildwheel}, the standard wheel-building tool in the scientific Python ecosystem, which manages the appropriate Docker containers for Linux builds and native runners for macOS and Windows automatically.

\subsubsection{GitHub Actions Workflow}
\label{sec:phase2-workflow}

The workflow \texttt{r2py\_kernsmooth/.github/workflows/python-publish.yml} is triggered on every GitHub Release publication. It consists of three jobs.

The \textbf{\texttt{build\_wheels}} job runs on the runner matrix \texttt{[ubuntu-latest,} \texttt{ubuntu-24.04-arm,} \texttt{macos-13,} \texttt{macos-latest,} \texttt{windows-latest]}. \texttt{ubuntu-latest} (Linux~x86\_64) and \texttt{ubuntu-24.04-arm} (Linux~aarch64, a native ARM64 runner that requires no QEMU emulation) each launch both \texttt{manylinux} (glibc) and \texttt{musllinux} (Alpine/musl-libc) Docker containers managed by cibuildwheel, producing PEP~600-compliant wheels for all configured Python versions on both Linux libc variants. \texttt{macos-13} and \texttt{macos-latest} build macOS~x86\_64 (Intel) and macOS~arm64 (Apple~Silicon) wheels natively, respectively. \texttt{windows-latest} builds Windows~x86\_64 wheels natively; a preceding step using the \texttt{msys2/setup-msys2} action installs \texttt{gfortran}, OpenBLAS, and \texttt{pkg-config} into the MSYS2 MinGW64 environment before cibuildwheel is invoked. All configuration options are read from the \texttt{[tool.cibuildwheel]} table in \texttt{pyproject.toml} (\S\ref{sec:phase2-cibw-config}).

The \textbf{\texttt{build\_sdist}} job runs on \texttt{ubuntu-latest}, executes \texttt{python~-m~build~--sdist}, and uploads the resulting \texttt{.tar.gz} archive. The sdist allows users on platforms with no pre-built wheel---and users with non-standard build environments---to compile from source.

The \textbf{\texttt{pypi-publish}} job, which depends on both preceding jobs, downloads all wheel and sdist artifacts into a single \texttt{dist/} directory and publishes them via \texttt{pypa/gh-action-pypi-publish@release/v1}. Authentication uses OIDC-based \emph{trusted publishing}: the job requests a short-lived identity token (\texttt{id-token:~write} permission) that PyPI validates against the registered GitHub repository and Actions environment (\texttt{environment:~name:~pypi}), eliminating the need to store any long-lived credential in GitHub~Secrets.

\subsubsection{Platform Coverage}
\label{sec:phase2-cibw-config}

The \texttt{[tool.cibuildwheel]} configuration in \texttt{pyproject.toml} is as follows:

\begin{lstlisting}[style=shell]
[build-system]
build-backend = "mesonpy"
requires = ["meson-python", "meson", "ninja", "numpy"]

[project]
name = "r2py_kernsmooth"
version = "0.1.0"
requires-python = ">=3.10"
dependencies = ["numpy", "scipy"]

[tool.cibuildwheel]
build = "cp310-* cp311-* cp312-* cp313-* cp314-*"
skip = ["*-manylinux_i686"]
build-frontend = "build"

[tool.cibuildwheel.linux]
before-all = "command -v dnf && dnf install -y openblas-devel || command -v yum && yum install -y openblas-devel || apk add --no-cache gfortran openblas-dev"

[tool.cibuildwheel.macos]
before-all = "brew install openblas"

[tool.cibuildwheel.macos.environment]
PKG_CONFIG_PATH = "$(brew --prefix openblas)/lib/pkgconfig"

[tool.cibuildwheel.windows]
before-build = "pip install delvewheel"
repair-wheel-command = "delvewheel repair -w {dest_dir} {wheel}"

[tool.cibuildwheel.windows.environment]
PATH = "C:\\msys64\\mingw64\\bin;{PATH}"
PKG_CONFIG_PATH = "C:\\msys64\\mingw64\\lib\\pkgconfig"
\end{lstlisting}

The principal design decisions are as follows.

\begin{description}
  \item[Python versions.] \texttt{cp310-*} through \texttt{cp314-*} cover CPython~3.10--3.14; 32-bit Linux (\texttt{*-manylinux\_i686}) is the only excluded target.

  \item[Build frontend.] \texttt{build-frontend~=~"build"} invokes the PEP~517 \texttt{build} tool rather than \texttt{pip~wheel}, which is more compatible with \texttt{meson-python}.

  \item[Linux dependencies.] \texttt{before-all} runs once per container environment (not once per Python version) and installs OpenBLAS via a package-manager detection chain: \texttt{dnf} for manylinux\_2\_28 / AlmaLinux~8 base images, \texttt{yum} for manylinux2014 / CentOS~7 base images, and \texttt{apk} for musllinux / Alpine base images. The \texttt{apk} branch additionally installs \texttt{gfortran}, since Alpine-based musllinux containers do not include it by default, unlike the manylinux containers.

  \item[macOS dependencies.] \texttt{brew install openblas} installs OpenBLAS via Homebrew. The \texttt{PKG\_CONFIG\_PATH} override directs Meson's \texttt{dependency('blas')} call to the Homebrew installation via \texttt{pkg-config}; without it, the first four stages of the BLAS discovery chain exhaust their search and the configure-time error in the fifth stage is raised.

  \item[Windows toolchain.] The \texttt{PATH} and \texttt{PKG\_CONFIG\_PATH} overrides expose the MSYS2 MinGW64 installation (set up by the preceding workflow step) to Meson inside each cibuildwheel build subprocess. \texttt{delvewheel}---the Windows counterpart of Linux \texttt{auditwheel} and macOS \texttt{delocate}---bundles the OpenBLAS DLL into the wheel, making it self-contained for users without MSYS2.
\end{description}

Across the five runners, up to 35 binary wheels are produced per release (five Python versions $\times$ seven platform/libc combinations: Linux~x86\_64 manylinux, Linux~x86\_64 musllinux, Linux~aarch64 manylinux, Linux~aarch64 musllinux, macOS~x86\_64, macOS~arm64, Windows~x86\_64).

\subsubsection{Coverage Limitations}

Three platform categories present in established scientific Python distributions such as scikit-learn are not covered.

\begin{description}
  \item[Windows arm64.] Native wheels for Qualcomm Snapdragon-based Windows PCs require a Fortran compiler for the Windows~arm64 ABI. The only readily available option is LLVM \texttt{flang} via the MSYS2 CLANGARM64 environment, but its integration with \texttt{f2py} and Meson is not yet production-stable. As an interim measure, Windows~arm64 users can install the \texttt{win\_amd64} wheel, which runs transparently under the Windows Prism x86-64 emulation layer.

  \item[Linux ppc64le and s390x.] IBM~POWER and IBM~Z mainframe architectures would require QEMU emulation in CI, at approximately 40~minutes per Python version per architecture. Given the extreme rarity of these platforms in scientific Python workflows, they are not included.

  \item[Free-threaded Python (cp313t, cp314t).] The experimental no-GIL CPython builds are supported by scikit-learn on selected platforms. Compatibility of the \texttt{f2py}-generated C extension with the free-threaded CPython ABI has not been verified and is deferred until the free-threaded build reaches stable status.
\end{description}

Before the first release, \texttt{r2py\_kernsmooth} must be registered on PyPI and the GitHub Actions \texttt{pypi} environment configured as a trusted publisher in the project's PyPI settings.

\section{Phase 3: Language Dependency Catalogue and Conversion Guides}
% =============================================================================

\subsection{Motivation}

The automated function conversion in Phase~4 maps each R construct to its Python equivalent. For most R standard-library functions, a na\"{i}ve translation to the apparent NumPy or SciPy analogue introduces silent numerical errors. Documented examples include: R's \texttt{fft(inverse=TRUE)} is unnormalised while NumPy's \texttt{ifft} normalises by $1/N$; R's \texttt{var()} uses Bessel's correction while \texttt{np.var} does not; R's \texttt{as.integer()} targets 32-bit integers while NumPy defaults to 64-bit.

One may argue that these discrepancies are well-known and can be handled on a case-by-case basis during the conversion process in Phase~4. However, due to the random nature of the LLM, we cannot guarantee that all language dependencies will be correctly identified and translated whenever a function is converted. If there is not a unified, machine-readable reference of all language dependencies and their translation nuances, then the risk of silent numerical errors in the final Python package is unacceptably high. (Even if the translation is mostly correct, we cannot ensure the consistency of choices or styles, which can make debugging significantly difficult.) Also, there are many language dependencies in R that have so many variations in their usage that it is hard to guarantee a correct translation without a comprehensive reference. For example, the \texttt{as.double()} function is used in many different contexts (such as on scalars, vectors, matrices, or even data frames), and its translation to Python can vary depending on the context. If we rely directly on the LLM to identify and translate each usage correctly during the conversion process, we may miss some cases or introduce errors. By contrast, if we systematically extract all language dependencies and their call sites in Phase~3, and then generate detailed conversion guides for each scenario for each dependency, we can ensure that the translation in Phase~4 is consistent and accurate across all functions.

Phase~3 systematically catalogs every R standard-library function used in \texttt{all.R}, records its precise call sites, and generates a dedicated translation guide for each. By encoding these nuances in machine-readable Markdown files that are later fed to the conversion sub-agent, the translation quality is improved without depending on the agent's training-data knowledge of R semantics solely.

\subsection{Language Dependency Extraction}

The \texttt{/extract-r-folder-language-dependencies} skill (Appendix~\ref{app:commands}, \S\ref{appentry:extract-r-folder-language-dependencies}) was invoked:

\begin{lstlisting}[style=shell]
/extract-r-folder-language-dependencies \
    KernSmooth/R/ \
    structural_analysis/R/
\end{lstlisting}

The skill dispatches the \texttt{@extract-r-file-language-dependencies} sub-agent (Appendix~\ref{app:agents}, \S\ref{appentry:extract-r-file-language-dependencies}) on each R file together with its pre-computed structural-analysis JSON. The sub-agent reads the structural JSON to identify the language dependencies of each function, then rescans the R source line by line to record, for every individual call site: the dependency name, the enclosing function name, the line number, and the complete call expression (including multi-line calls).

The result was written to \texttt{language\_dependency\_analysis/R/all.csv}: a 370-row, 4-column CSV (\texttt{language\_dependency}, \texttt{function\_name}, \texttt{line\_number}, \texttt{call\_body}) covering all 16 functions in \texttt{all.R}.

A utility script \texttt{language\_dependency\_analysis/combine\_csvs.py} was also written to:
\begin{itemize}
  \item Recursively glob all \texttt{.csv} files under \texttt{language\_dependency\_analysis/R/}.
  \item Insert a \texttt{file\_path} column (the relative path of the source R file within the \texttt{R/} folder, \texttt{.csv} extension replaced with \texttt{.R}).
  \item Concatenate, sort by \texttt{[language\_dependency, file\_path, function\_name, line\_number]}, and save the combined result to \texttt{language\_dependency\_analysis/combined\_table.csv}.
\end{itemize}

This script is designed to sort the entries in the original \texttt{all.csv} to make the entries for the same language dependency contiguous and convenient to be extracted for further tasks. It also makes it easier to support future analysis of additional R files beyond \texttt{all.R}.

\subsection{Conversion Guide Generation}

The \texttt{/generate-language-dependency-conversion-guides} skill (Appendix~\ref{app:commands}, \S\ref{appentry:generate-language-dependency-conversion-guides}) was invoked:

\begin{lstlisting}[style=shell]
/generate-language-dependency-conversion-guides \
    KernSmooth/R/ \
    language_dependency_analysis/combined_table.csv
\end{lstlisting}

The skill identified 46 unique \texttt{language\_dependency} values in the combined CSV and launched 46 instances of the \texttt{@generate-language-dependency-conversion-guide} sub-agent (Appendix~\ref{app:agents}, \S\ref{appentry:generate-language-dependency-conversion-guide}) in a single parallel batch. Each agent:
\begin{enumerate}
  \item Received the CSV subset for its assigned dependency.
  \item Read the corresponding source lines in \texttt{all.R} for contextual understanding of argument types and surrounding logic.
  \item Consulted R online documentation (\url{https://www.rdocumentation.org}) where the function's behavior was complex or non-obvious.
  \item Produced a structured Markdown guide specifying the correct Python translation strategy, with concrete code examples.
\end{enumerate}

The 46 resulting guides were saved to \texttt{language\_dependency\_analysis/conversion\_guides/} with filenames matching the dependency name (e.g.\ \texttt{fft.md}, \texttt{var.md}, \texttt{match.arg.md}).

\subsection{Selected Translation Nuances}

The following are the most consequential findings from the conversion guides, as they directly informed specific decisions in Phase~4. Together, they constitute the principal risk of silent numerical error in any R-to-Python translation of this package.

\subsubsection{FFT Normalisation (\texttt{fft.md})}

R's \texttt{fft(\ldots,\,inverse=TRUE)} returns the \emph{raw, unnormalised} inverse DFT: for an $N$-point transform it returns $\sum_{k=0}^{N-1} X_k e^{2\pi i k n /N}$ without dividing by $N$. NumPy's \texttt{np.fft.ifft} normalises by $1/N$ by default, as per the conventional definition. The R source explicitly divides the inverse-FFT result by the grid size $P$ (or $P_1 \times P_2$ for two-dimensional transforms) after computing \texttt{Re()}. When porting to NumPy, those division operations must be \emph{removed}.

\subsubsection{Fortran FFI Type Coercions (\texttt{as.double.md}, \texttt{as.integer.md})}

R's pre-allocation idiom \texttt{as.double(rep(0, n))} maps to \texttt{np.zeros(n, dtype=np.float64)}. \texttt{as.integer(Lvec)} on a vector maps to \texttt{Lvec.astype(np.int32)}. The 32-bit integer constraint is non-obvious but critical: standard Fortran \texttt{INTEGER} is 32~bits; passing \texttt{np.int64} arrays would silently corrupt the values seen by the Fortran subroutine on a 64-bit system.

\subsubsection{Sample Variance Denominator (\texttt{var.md})}

R's \texttt{var(x)} uses Bessel's correction ($n-1$ denominator), i.e.\ it computes the \emph{sample} variance. NumPy's \texttt{np.var} defaults to the \emph{population} variance ($n$ denominator, \texttt{ddof=0}). All occurrences of \texttt{sqrt(var(x))} in the R source must be translated as \texttt{np.std(x, ddof=1)}.

\subsubsection{Optional Argument Detection (\texttt{missing.md})}

R's \texttt{missing(param)} returns \texttt{TRUE} if the caller did not supply the argument. Python has no built-in equivalent. The idiomatic replacement is the \texttt{None} sentinel default: \texttt{def f(x, param=None)} with \texttt{if param is None:} guards inside the function body.

\subsubsection{Negative-Base Fractional Exponentiation (\texttt{sqrt.md})}

Several bandwidth formulas involve odd roots of quantities that may be negative, e.g.\ $(-3\sqrt{2/\pi}\,/\,\hat{\psi}_6)^{1/7}$. R silently computes real-valued odd roots of negative numbers via its internal rule for fractional exponents. Python's \texttt{**} operator raises \texttt{ValueError} for negative bases with non-integer exponents. The correct idiom is
\begin{lstlisting}[style=python]
math.copysign(abs(x) ** (1.0 / n), x)
\end{lstlisting}

\subsubsection{Array Memory Order (\texttt{matrix.md})}

R's \texttt{matrix()} and Fortran both use \emph{column-major} (Fortran-order) storage. NumPy defaults to \emph{row-major} (C-order). When reshaping arrays produced by Fortran subroutines, \texttt{np.reshape(\ldots,\,order=`F')} is required.

\subsubsection{Package Lifecycle Hooks (\texttt{library.dynam.unload.md})}

R provides package attach/detach lifecycle hooks via \texttt{.onAttach} and \texttt{.onUnload}. \texttt{library.dynam.unload} has no Python equivalent: CPython does not safely support unloading native extension modules. The guide documents that the \texttt{.onUnload} body should either be omitted or replaced with an \texttt{atexit} callback for non-library cleanup tasks.

\subsubsection{Partial String Matching (\texttt{match.arg.md})}

R's \texttt{match.arg(arg, choices)} validates \texttt{arg} against a vector of allowed strings and supports \emph{unambiguous prefix matching} (e.g.\ \texttt{"ep"} matches \texttt{"epanechnikov"}). No standard Python equivalent exists. The guide specifies a reusable helper function that performs exact match, then unique-prefix match, then raises a descriptive error.

\subsubsection{R \texttt{switch} on Strings (\texttt{switch.md})}

R's \texttt{switch(key, case1=expr1, case2=expr2, \ldots)} on string keys translates cleanly to a Python dictionary lookup. For cheap scalar computations (as in \texttt{KernSmooth}), a plain \texttt{dict} of values is sufficient; a \texttt{dict} of \texttt{lambda} functions is available for cases requiring lazy evaluation.

% =============================================================================
\section{Phase 4: Automated Function Conversion and Package Assembly}
% =============================================================================

\subsection{Motivation}

With the dependency map (Phase~1), build infrastructure (Phase~2), and translation guides (Phase~3) in place, Phase~4 executes the actual conversion. Two challenges make this non-trivial even with the preparation work done:

\begin{enumerate}
  \item Functions must be converted in dependency order: when the conversion agent for a function reads the JSON files of its callees to understand their Python calling conventions, those files must already exist.
  \item The Fortran FFI calling convention in Python (via \texttt{f2py}) differs fundamentally from R's \texttt{.Fortran()} in how it handles output arguments. This difference is invisible from the R source and therefore not naturally handled by pattern matching; it requires targeted post-conversion patching.
\end{enumerate}

\subsection{Function Conversion}

The \texttt{/convert-r-file-to-python} skill (Appendix~\ref{app:commands}, \S\ref{appentry:convert-r-file-to-python}) was invoked:

\begin{lstlisting}[style=shell]
/convert-r-file-to-python \
    KernSmooth/R/all.R \
    structural_analysis/R/all.json \
    structural_analysis/dependency_levels.csv \
    language_dependency_analysis/conversion_guides/ \
    conversion_results/R/
\end{lstlisting}

The skill parses \texttt{dependency\_levels.csv}, filters for \texttt{r\_file == all.R}, and derives a topological ordering from leaves to roots. It then invokes the \texttt{@convert-r-function-to-python} sub-agent (Appendix~\ref{app:agents}, \S\ref{appentry:convert-r-function-to-python}) once per function, \emph{sequentially} in topological order, providing each agent with:
\begin{itemize}
  \item The R source fragment for the function.
  \item The function's JSON dependency map (from \texttt{all.json}).
  \item The folder of 46 language-dependency conversion guides.
  \item The output folder \texttt{conversion\_results/R/all.R/} for the resulting JSON artefact.
\end{itemize}

The conversion order was:

\begin{center}
\begin{tabular}{ll}
\toprule
Level & Functions (in conversion order) \\
\midrule
2 (deepest leaves) & \texttt{linbin}, \texttt{rlbin} \\
1 (intermediate) &
  \texttt{bkfe}, \texttt{blkest}, \texttt{cpblock},
  \texttt{linbin2D}, \texttt{locpoly}, \texttt{sdiag}, \texttt{sstdiag} \\
0 (roots) &
  \texttt{.onAttach}, \texttt{.onUnload}, \texttt{bkde},
  \texttt{bkde2D}, \texttt{dpih}, \texttt{dpik}, \texttt{dpill} \\
\bottomrule
\end{tabular}
\end{center}

The 16 resulting JSON artefacts were written to \texttt{conversion\_results/R/all.R/}. (Two artefacts --- \texttt{.onAttach.json} and \texttt{.onUnload.json} --- use leading-dot filenames and are not visible to standard \texttt{ls}; they require \texttt{ls -a} to discover.)

\subsubsection{Notable Per-Function Translation Decisions}

\begin{description}
  \item[\texttt{linbin}/\texttt{rlbin}/\texttt{linbin2D}/\texttt{blkest}/\texttt{cpblock}]
    \texttt{.Fortran(F\_xxx,~\ldots)[[k]]} mapped to \texttt{\_KernSmooth.xxx(\ldots)} with pre-allocated NumPy output buffers.
  \item[\texttt{bkfe}]
    Hermite polynomial recurrence loop translated directly; the manual \texttt{/P} division present in the R source was removed per the \texttt{fft.md} guide, since NumPy's \texttt{np.fft.ifft} already normalises.
  \item[\texttt{locpoly}]
    Dual density/regression branch (controlled by whether the \texttt{y} argument is \texttt{None}); a variable-length bandwidth discretization loop; a 19-argument Fortran subroutine call; and a \texttt{math.gamma(drv+1)} post-multiplier for the derivative-order correction.
  \item[\texttt{sdiag}/\texttt{sstdiag}]
    The \texttt{ss}, \texttt{Smat}, \texttt{uu}, \texttt{Umat} matrices are flattened column-major (\texttt{order=`F'}) before being passed to the Fortran subroutine, consistent with the \texttt{matrix.md} guide.
  \item[\texttt{dpih}/\texttt{dpik}]
    Negative-base fractional exponents (e.g.\ $(-3\sqrt{2/\pi}/\hat\psi_6)^{1/7}$) replaced with \texttt{math.copysign(abs(val)**(1/k), val)} per the \texttt{sqrt.md} guide.
  \item[\texttt{bkde2D}]
    Two-dimensional FFT convolution via \texttt{np.fft.fft2}/\texttt{np.fft.ifft2}; wrap-around kernel construction using Python slice notation to replicate R's decreasing-index sequences.
  \item[\texttt{.onAttach}]
    \texttt{packageStartupMessage} translated to \texttt{print(\ldots,\,file=sys.stderr)}.
  \item[\texttt{.onUnload}]
    \texttt{library.dynam.unload} body omitted per \texttt{library.dynam.unload.md}; function registered as an \texttt{atexit} callback at module level.
\end{description}

\subsection{Package Assembly}

The \texttt{/combine-python-functions-into-file} skill (Appendix~\ref{app:commands}, \S\ref{appentry:combine-python-functions-into-file}) was invoked:

\begin{lstlisting}[style=shell]
/combine-python-functions-into-file \
    conversion_results/R/all.R/ \
    r2py_kernsmooth/r2py_kernsmooth/__init__.py
\end{lstlisting}

Before writing, the skill performed a structural scan of \texttt{meson.build} and the existing \texttt{\_\_init\_\_.py} to verify the correct import path for the Fortran extension (\texttt{from~.~import~\_KernSmooth}), then corrected discrepancies found in the per-function JSON artefacts:

\begin{center}
\begin{tabular}{p{0.44\textwidth}p{0.44\textwidth}}
\toprule
Original (in JSON artefacts) & Corrected \\
\midrule
\texttt{from r2py\_kernsmooth import \_KernSmooth} &
  \texttt{from . import \_KernSmooth} \\
\texttt{from r2py\_kernsmooth.linbin import linbin} &
  removed (functions are co-located) \\
\texttt{from scipy.stats import norm} (duplicate) &
  merged: \texttt{from scipy.stats import beta as beta\_dist, norm} \\
\bottomrule
\end{tabular}
\end{center}

The final import block, sorted standard library $\to$ third-party $\to$ local, was:

\begin{lstlisting}[style=python]
import atexit
import math
import sys
import warnings

import numpy as np
from scipy.stats import beta as beta_dist, norm

from . import _KernSmooth
\end{lstlisting}

The \texttt{\_\_all\_\_} list was populated with all 14 public function names (excluding \texttt{\_on\_attach} and \texttt{\_on\_unload}). Functions were written to the file in leaf-to-root order, ensuring all callees are defined before callers. Syntax was confirmed immediately by \texttt{ast.parse()}.

\subsection{The f2py Return-Value Bug}

\subsubsection{Observed Symptom}

After assembly, running the test script raised:

\begin{lstlisting}[style=shell]
TypeError: 'NoneType' object is not subscriptable
  File "__init__.py", line 36, in linbin
    gcnts = _KernSmooth.linbin(...)[6]
\end{lstlisting}

\subsubsection{Root Cause: R \texttt{.Fortran()} vs.\ f2py Semantics}

R's \texttt{.Fortran()} always returns a \emph{named list of all arguments}, including modified output arrays. The R idiom \texttt{.Fortran(F\_linbin, \ldots)\$gcnts} (or equivalently \texttt{[[6]]}) extracts the modified output array from that list.

f2py-wrapped Fortran \emph{subroutines} return \texttt{None} unconditionally. Output array arguments are modified in-place via C pointer; the Python caller reads results by examining the pre-allocated arrays that were passed in. Scalar output arguments require special treatment: passing a Python \texttt{np.float64(0.0)} (immutable) does not allow Fortran to write back. The fix is to use a length-1 NumPy array (\texttt{np.zeros(1, dtype=np.float64)}), which is mutable and passes a writable pointer.

The \texttt{@convert-r-function-to-python} agent for \texttt{linbin} had generated \texttt{[\,6\,]} subscripting on the \texttt{None} return. The agent for \texttt{rlbin} (converted first) had correctly identified the in-place semantics and avoided this. The inconsistency arose because each agent operated independently.

\subsubsection{Fortran Output Argument Audit}

All Fortran source files were read to identify output arguments:

\begin{center}
\begin{tabular}{lll}
\toprule
Subroutine & Output argument(s) & Type \\
\midrule
\texttt{linbin} & \texttt{gcnts(*)} & \texttt{double precision} array \\
\texttt{lbtwod} & \texttt{gcnts(*)} & \texttt{double precision} array \\
\texttt{blkest} & \texttt{sigsqe}, \texttt{th22e}, \texttt{th24e} &
  \texttt{double precision} \emph{scalars} \\
\texttt{cp}     & \texttt{Cpvals(Nmax)} & \texttt{double precision} array \\
\texttt{locpol} & \texttt{cvest(*)} & \texttt{double precision} array \\
\texttt{sdiag}  & \texttt{Sdg(*)} & \texttt{double precision} array \\
\texttt{sstdg}  & \texttt{SSTd(*)} & \texttt{double precision} array \\
\bottomrule
\end{tabular}
\end{center}

\subsubsection{Patches Applied}

Seven targeted edits were made to \texttt{r2py\_kernsmooth/r2py\_kernsmooth/\_\_init\_\_.py}:

\begin{enumerate}
  \item \textbf{\texttt{linbin}} --- Removed \texttt{[6]} subscript; read \texttt{gcnts} (pre-allocated \texttt{np.zeros(M)}) directly after the call.
  \item \textbf{\texttt{linbin2D}} --- Removed \texttt{out = \ldots} and \texttt{out[8]}; read \texttt{gcnts} then reshaped with \texttt{order=`F'}.
  \item \textbf{\texttt{blkest}} --- Changed \texttt{sigsqe/th22e/th24e} from \texttt{np.float64(0.0)} to \texttt{np.zeros(1, dtype=np.float64)}; return values read as \texttt{sigsqe[0]}, \texttt{th22e[0]}, \texttt{th24e[0]}.
  \item \textbf{\texttt{cpblock}} --- Removed \texttt{out = \ldots} and \texttt{out[12]}; used \texttt{Cpvals} (modified in-place) directly in \texttt{np.argmin}.
  \item \textbf{\texttt{locpoly}} --- Removed \texttt{out = \ldots} and \texttt{out[18]}; applied \texttt{math.gamma(drv+1)} multiplier to \texttt{curvest} directly.
  \item \textbf{\texttt{sdiag}} --- Removed \texttt{out = \ldots} and \texttt{out[16]}; returned \texttt{Sdg} directly.
  \item \textbf{\texttt{sstdiag}} --- Removed \texttt{out = \ldots} and \texttt{out[18]}; returned \texttt{SSTd} directly.
\end{enumerate}

A post-patch \texttt{grep} confirmed zero remaining stale \texttt{out[integer]} subscripts on Fortran call returns. \texttt{ast.parse()} re-confirmed syntax validity.

\subsection{Output Validation}

The test script \texttt{r2py\_kernsmooth/tests/test.py} invoked \texttt{r2py\_kernsmooth.bkde(np.array([1,2,3,4,5]))} with default parameters and printed the result. The output was verified analytically:

\begin{itemize}
  \item \textbf{Grid range:} $[-4.244,\,10.244]$. Expected $[\min(x)-\tau h,\,\max(x)+\tau h]$ where $\tau=4$ and $h \approx \delta_0 \cdot (243/175)^{1/5} \cdot \sigma \approx 0.776 \times 1.069 \times 1.581 \approx 1.312$, giving $[-4.248,\,10.248]$. \checkmark
  \item \textbf{Grid size:} 401 points (default \texttt{gridsize=401}).~\checkmark
  \item \textbf{Peak location:} Maximum density $\approx 0.18989$ at $x=3.0$, matching the mean and median of the symmetric input $\{1,2,3,4,5\}$.~\checkmark
  \item \textbf{Symmetry:} The density array is a palindrome.~\checkmark
  \item \textbf{Non-negativity:} All values $>0$.~\checkmark
  \item \textbf{Normalisation:} Trapezoidal integral $\approx 1.0$.~\checkmark
\end{itemize}

After Phase~4, the package was in a runnable state with all 16 functions converted.

% =============================================================================
\section{Phase 5: Code Quality Audit and Namespace Refinement}
% =============================================================================

\subsection{Motivation}

The automated conversion pipeline prioritises functional correctness. Phase~5 is a manual audit and refactoring pass that addresses three categories of follow-on concerns:

\begin{enumerate}
  \item \textbf{Type annotations:} The conversion agent was instructed to produce complete annotations, but the initial assembly contained bare container types (\texttt{np.ndarray}, \texttt{dict}) without subtype parameters.
  \item \textbf{Namespace hygiene:} Module-level imports (\texttt{import numpy as np}, etc.) are visible as package attributes; best-practice conventions for managing this were investigated and applied.
  \item \textbf{Code quality:} A systematic audit identified dead code left over from the R porting process, a latent crash bug, structural duplication, and PEP~8 violations, all of which were resolved.
\end{enumerate}

\subsection{Type Annotation Completion}

All 14 public function signatures and both private function signatures were reviewed. The following annotation deficiencies were found and corrected:

\begin{itemize}
  \item Bare \texttt{np.ndarray} parameters and return types, which lack dtype information.
  \item Bare \texttt{dict}, \texttt{tuple}, and \texttt{list} without element type parameters.
  \item Optional parameters with default \texttt{None} missing the \texttt{| None} union form (e.g.\ \texttt{bandwidth}, \texttt{range\_x}, \texttt{y}, \texttt{degree}).
\end{itemize}

\texttt{from typing import Any} was added to the import block. Every array annotation was expanded to \texttt{np.ndarray[Any, np.dtype[np.float64]]}, reflecting the exclusive use of 64-bit floating-point arrays throughout the module. Return types of functions returning dictionaries were updated:

\begin{itemize}
  \item \texttt{rlbin}, \texttt{locpoly}, \texttt{sdiag}, \texttt{sstdiag}, \texttt{bkde}, \texttt{bkde2D} $\to$ \texttt{dict[str, np.ndarray[Any, np.dtype[np.float64]]]}
  \item \texttt{blkest} $\to$ \texttt{dict[str, np.float64]} (scalar values extracted from length-1 arrays)
\end{itemize}

The \texttt{NDArray[np.float64]} alias from \texttt{numpy.typing} was evaluated and rejected in favour of the explicit \texttt{np.ndarray} base type.

\subsection{Namespace Privacy Investigation}

Importing \texttt{numpy as np} at module level exposes \texttt{r2py\_kernsmooth.np} as a public package attribute, which is undesirable but difficult to prevent. Two strategies were evaluated:

\begin{description}
  \item[Underscore aliasing] All imports renamed to \texttt{\_np}, \texttt{\_math}, etc., with global replacement of all references throughout the file. Implemented and tested (confirmed \texttt{hasattr(r2py\_kernsmooth, `np')} returned \texttt{False}), then \textbf{reverted}. A survey of major scientific Python packages (NumPy, SciPy, pandas) confirmed that underscore-aliasing of dependency imports is not standard practice; all use conventional import names at module level.
  \item[\texttt{del} cleanup] Applicable only to imports used \emph{exclusively at module initialisation time}. Only \texttt{atexit} qualified: it appeared solely in the \texttt{atexit.register(\_on\_unload)} call. \texttt{del atexit} was added immediately after that registration. No other import could be safely deleted because all others (\texttt{math}, \texttt{np}, etc.) are referenced inside function bodies, where deletion would cause \texttt{NameError} at call time.
\end{description}

PEP~562 (\texttt{module.\_\_getattr\_\_}) was also evaluated but found inapplicable: it is invoked only for \emph{missing} attributes and cannot hide names already present in \texttt{\_\_dict\_\_}. The \texttt{\_\_all\_\_} list was confirmed as the authoritative and conventional public-API contract.

\subsection{Full Code Audit and Systematic Refactoring}

A line-by-line review of \texttt{\_\_init\_\_.py} identified seven defect categories. All were resolved in a single comprehensive rewrite.

\subsubsection{Dead Code from R Porting}

\texttt{\_on\_attach} and \texttt{\_on\_unload} were direct transliterations of R's \texttt{.onAttach} and \texttt{.onUnload} lifecycle hooks. Python's import machinery provides no equivalent of R's package attach/detach events; neither function was reachable from any Python code path. Their presence also rendered \texttt{import atexit}, \texttt{import sys}, and the subsequently added \texttt{del atexit} statement all useless. All were removed.

An R-era source-control comment inside \texttt{cpblock} (\texttt{\# remove unused `q' 2007-07-10}) was also removed; it had no meaning in a Python file and referenced an event in the R package's version history.

\subsubsection{Silent Crash Bug in Three Public Functions}

\texttt{locpoly}, \texttt{sdiag}, and \texttt{sstdiag} all declared \texttt{bandwidth:~np.ndarray[\ldots]\,|\,None~=~None} (a \texttt{None} default with a \texttt{| None} annotation added during type annotation completion) yet called \texttt{np.asarray(bandwidth, dtype=np.float64)} unconditionally. Passing \texttt{None} produced a 0-d array with value \texttt{nan}; a subsequent \texttt{len(bandwidth)} call on that 0-d array raised \texttt{TypeError:~len()~of~unsized~object}.

The annotation advertised \texttt{None} as a semantically valid input while the implementation crashed silently. In the internal call graph, the bug was unreachable because all three functions are always called from \texttt{dpill} with a concrete bandwidth value --- but any external caller relying on the type annotation would have encountered it immediately. The resolution mirrors R's \texttt{!missing(bandwidth)} semantics via the \texttt{None}-sentinel idiom from the \texttt{missing.md} conversion guide: the bandwidth positivity check in \texttt{locpoly} was corrected from an unconditional form to:

\begin{lstlisting}[style=python]
if bandwidth is not None and np.any(np.asarray(bandwidth) <= 0):
    raise ValueError("'bandwidth' must be strictly positive")
\end{lstlisting}

This matches R's \texttt{if (!missing(bandwidth) \&\& any(bandwidth <= 0))} guard. \texttt{sdiag} and \texttt{sstdiag} carry no bandwidth positivity check in the R source and required no guard.

\subsubsection{Structural Code Duplication}

Two logic blocks appeared identically across three functions each:

\begin{description}
  \item[Bandwidth discretization (27 lines)] In \texttt{locpoly}, \texttt{sdiag}, and \texttt{sstdiag}: handles the scalar-bandwidth path, the vector-of-length-$M$ path, and an error path. Extracted into:
    \begin{lstlisting}[style=python]
def _discretize_bandwidth(bandwidth, M, delta, Q, tau):
    ...  # returns (hdisc, Lvec, indic, Q)
    \end{lstlisting}
    The extraction also corrected a robustness defect in the scalar-detection logic. The original code used \texttt{elif len(bandwidth) == 1}, which raises \texttt{TypeError} for 0-d NumPy arrays produced by \texttt{np.asarray} on a plain Python \texttt{float} or \texttt{np.float64} scalar. The corrected guard, consistent with the pattern already used in \texttt{bkde2D}, is:
    \begin{lstlisting}[style=python]
if bw.ndim == 0 or len(bw) == 1:
    ...
    \end{lstlisting}
  \item[Prefix-match validation (8 lines)] In \texttt{bkde} (for \texttt{kernel}), \texttt{dpih} (for \texttt{scalest}), and \texttt{dpik} (for both \texttt{kernel} and \texttt{scalest}). Extracted into:
    \begin{lstlisting}[style=python]
def _resolve_choice(val, choices):
    ...  # exact match -> unique prefix match -> error
    \end{lstlisting}
    The error message format was subsequently corrected to match R's \texttt{match.arg} output; see Section~\ref{sec:phase5-r-alignment}.
\end{description}

Both helpers are private (underscore-prefixed) and absent from \texttt{\_\_all\_\_}, preserving the public API surface.

\subsubsection{PEP~8 Signature Reformatting}

All function signatures exceeding 88 characters were reformatted to multi-line style with one parameter per line. The most extreme cases (\texttt{locpoly}, \texttt{sdiag}, \texttt{sstdiag}) had reached 210+ characters on a single line prior to reformatting.

\subsubsection{Error Handling Alignment with the R Source}
\label{sec:phase5-r-alignment}

A subsequent audit compared every \texttt{stop()}, \texttt{warning()}, and \texttt{match.arg()} call in \texttt{KernSmooth/R/all.R} against the corresponding \texttt{raise}/\texttt{warn} in \texttt{\_\_init\_\_.py}. Three discrepancies were identified and corrected.

\begin{description}
  \item[\texttt{\_resolve\_choice} error message.]
    R's \texttt{match.arg} produces a single message for both ambiguous partial matches and non-matching values, irrespective of which failure mode triggered:
    \begin{lstlisting}[style=shell]
'arg' should be one of "c1", "c2", ...
    \end{lstlisting}
    The original \texttt{\_resolve\_choice} generated two distinct custom-formatted messages, neither matching this format. Both failure paths were unified to produce the R-equivalent message, and the now-unused \texttt{param\_name} parameter was removed from the function signature and all four call sites.

  \item[Invented \texttt{None} guard in \texttt{locpoly}, \texttt{sdiag}, \texttt{sstdiag}.]
    An intermediate draft had added an explicit \texttt{ValueError("'bandwidth'~must~be~specified")} guard at the top of each function body. No corresponding \texttt{stop()} call exists in the R source: when a required argument is absent, R's runtime raises ``argument is missing, with no default'' at the language level rather than via a user-defined message. The guard was removed from all three functions to eliminate error behavior with no R counterpart.

  \item[\texttt{locpoly} positivity check pattern.]
    R guards its bandwidth positivity check with \texttt{!missing(bandwidth)}: \texttt{if (!missing(bandwidth) \&\& any(bandwidth <= 0))}. The Python positivity check in \texttt{locpoly} had been written in unconditional form, diverging from this pattern. It was corrected to apply the \texttt{None}-sentinel idiom (\texttt{bandwidth is not None and \ldots}), as documented in the \texttt{missing.md} conversion guide.
\end{description}

All remaining \texttt{stop()} and \texttt{warning()} messages in the R source were confirmed as exactly reproduced in the Python module.

\subsection{State of the Package After Phase~5}

The rewritten \texttt{r2py\_kernsmooth/r2py\_kernsmooth/\_\_init\_\_.py} satisfies the following properties:

\begin{itemize}
  \item No dead code or unused imports.
  \item No bare \texttt{np.ndarray}, \texttt{dict}, \texttt{tuple}, or \texttt{list} annotations in any public or private function prototype.
  \item All error and warning messages exactly match the explicit \texttt{stop()} and \texttt{warning()} calls in \texttt{KernSmooth/R/all.R}; no invented error messages are present.
  \item The \texttt{None}-sentinel idiom (\texttt{bandwidth is not None and \ldots}) consistently mirrors R's \texttt{!missing(bandwidth)} predicate across all functions that test bandwidth before use.
  \item No duplicated bandwidth-discretization or argument-validation logic.
  \item Correct 0-d NumPy array handling in the scalar-bandwidth code path.
  \item PEP~8 line-length compliance ($\leq 88$ characters) across all function signatures.
  \item No R-era changelog comments or R-specific lifecycle hook documentation.
\end{itemize}

The runtime target is CPython~3.14 (confirmed by the presence of \texttt{\_\_pycache\_\_/\_\_init\_\_.cpython-314.pyc}). The package uses Python~3.10+ syntax (\texttt{dict[str,~\ldots]}, \texttt{float\,|\,None} union types, lowercase generics), which is consistent with the runtime version.

% =============================================================================
% Bibliography
% =============================================================================

\begin{thebibliography}{9}
\bibitem{wandjones1995}
  M.~P.~Wand and M.~C.~Jones (1995).
  \textit{Kernel Smoothing}.
  Chapman and Hall, London.

\bibitem{sheatherjones1991}
  S.~J.~Sheather and M.~C.~Jones (1991).
  A reliable data-based bandwidth selection method for kernel density estimation.
  \textit{Journal of the Royal Statistical Society, Series B}, \textbf{53}(3), 683--690.

\bibitem{ruppertsheatherwand1995}
  D.~Ruppert, S.~J.~Sheather, and M.~P.~Wand (1995).
  An effective bandwidth selector for local least squares regression.
  \textit{Journal of the American Statistical Association}, \textbf{90}(432), 1257--1270.

\bibitem{wand1994}
  M.~P.~Wand (1994).
  Fast computation of multivariate kernel estimators.
  \textit{Journal of Computational and Graphical Statistics}, \textbf{3}(4), 433--445.
\end{thebibliography}

% =============================================================================
% Appendices
% =============================================================================

\clearpage
\begin{appendices}

% -----------------------------------------------------------
\section{Agent Specifications}
\label{app:agents}
% -----------------------------------------------------------

This appendix reproduces the complete specifications of the sub-agents deployed in this project. Each file is stored under \texttt{.claude/agents/} in the project repository. The YAML front matter names the agent and provides a one-line description; the body specifies the execution steps and output schema in structured natural language. These files are the normative definitions of agent behaviour: what the agent receives as input, how it processes it, and what it must output.

\subsection{Agents Used in Phase~1}

\appendixentry{analyze-r-file-dependencies}{../../.claude/agents/analyze-r-file-dependencies.md}

\subsection{Agents Used in Phase~3}

\appendixentry{extract-r-file-language-dependencies}{../../.claude/agents/extract-r-file-language-dependencies.md}

\appendixentry{generate-language-dependency-conversion-guide}{../../.claude/agents/generate-language-dependency-conversion-guide.md}

\subsection{Agents Used in Phase~4}

\appendixentry{convert-r-function-to-python}{../../.claude/agents/convert-r-function-to-python.md}

\subsection{Supplementary Agent}

The following agent is defined in the repository but has not yet been invoked in this project. It is included for completeness and as a reference for the Fortran source analysis task identified as a potential next step in Phase~1.

\appendixentry{analyze-c-file-dependencies}{../../.claude/agents/analyze-c-file-dependencies.md}

% -----------------------------------------------------------
\section{Skill (Command) Specifications}
\label{app:commands}
% -----------------------------------------------------------

This appendix reproduces the complete specifications of the top-level skills (slash commands) deployed in this project. Each file is stored under \texttt{.claude/commands/}. Skills are the entry points invoked interactively; they orchestrate one or more sub-agents (Appendix~\ref{app:agents}) and handle batch dispatch, error recovery, and output persistence.

\subsection{Skills Used in Phase~1}

\appendixentry{analyze-r-folder-dependencies}{../../.claude/commands/analyze-r-folder-dependencies.md}

\subsection{Skills Used in Phase~3}

\appendixentry{extract-r-folder-language-dependencies}{../../.claude/commands/extract-r-folder-language-dependencies.md}

\appendixentry{generate-language-dependency-conversion-guides}{../../.claude/commands/generate-language-dependency-conversion-guides.md}

\subsection{Skills Used in Phase~4}

\appendixentry{convert-r-file-to-python}{../../.claude/commands/convert-r-file-to-python.md}

\appendixentry{combine-python-functions-into-file}{../../.claude/commands/combine-python-functions-into-file.md}

\subsection{Workflow Utility}

The following skill was used at the end of each phase to produce the structured phase summaries stored in \texttt{docs/daily\_summaries/}, which are the primary source material for this report.

\appendixentry{summarize-research-report}{../../.claude/commands/summarize-research-report.md}

\end{appendices}

\end{document}
